\section{Công trình liên quan (Related Work)}
\label{sec:related}

Chúng tôi nhóm related work theo bốn trục khớp với pipeline: (i) truy hồi claim khoa học và SciFact;
(ii) sparse, dense, hybrid và các biến thể learned-sparse / late-interaction; (iii) routing và
cost-aware inference; (iv) fine-tuning hiệu quả, reranking, query performance prediction, và conformal
risk control.

\subsection{Truy hồi claim khoa học và SciFact}
SciFact giới thiệu bài toán xác minh claim khoa học: chọn các research abstract chứa bằng chứng cho một
claim, rồi xác định rationale và nhãn như SUPPORTS hoặc REFUTES~\citep{wadden2020scifact}. Benchmark
\textbf{BEIR} sau đó đóng gói SciFact như một phần của một benchmark zero-shot retrieval đa miền, thuận
tiện cho việc đánh giá retriever dưới một interface truy hồi nhất quán~\citep{thakur2021beir}. Dự án này
áp dụng setup SciFact theo phong cách BEIR và đánh giá việc một claim có truy hồi đúng các abstract liên
quan của nó hay không.

\subsection{Sparse, Dense và Hybrid retrieval}
BM25 là một hàm xếp hạng lexical cổ điển xuất phát từ probabilistic relevance framework~\citep{robertson2009bm25};
nó vẫn là một baseline mạnh cho tìm kiếm chuyên ngành vì exact term matching thường quan trọng. Dense
retrieval biểu diễn query và document thành vector và xếp hạng theo embedding similarity. Trong miền khoa
học, \textbf{SPECTER} dùng citation-informed transformer pretraining cho scientific document
embeddings~\citep{cohan2020specter}, còn \textbf{SciNCL} cải thiện biểu diễn tài liệu khoa học bằng
neighborhood contrastive learning trên thông tin citation graph~\citep{ostendorff2022scincl}. Hybrid
retrieval kết hợp sparse và dense; dự án dùng \textbf{Reciprocal Rank Fusion} (RRF), một phương pháp rank
fusion đơn giản nhưng hiệu quả để kết hợp các hệ thống retrieval khác nhau~\citep{cormack2009rrf}.

\subsection{Learned-sparse và Late-interaction}
\textbf{SPLADE} học các biểu diễn sparse lexical expansion vẫn tương thích với inverted index trong khi
cải thiện chất lượng neural retrieval~\citep{formal2021splade}. \textbf{ColBERT} dùng contextualized late
interaction để có matching mạnh hơn trong khi precompute biểu diễn document offline~\citep{khattab2020colbert}.
Đây là các phương án thay thế quan trọng, nhưng \emph{chưa} được hiện thực trong phase hiện tại vì dự án
giới hạn trong khung 3 tháng và phần cứng vừa phải; chúng được thảo luận như related work và hướng mở rộng.

\subsection{Routing và Cost-aware inference}
Các phương pháp routing chọn giữa các hệ thống thay thế dựa trên input. Trong bối cảnh LLM, \textbf{RouteLLM}
nghiên cứu routing giữa mô hình mạnh và yếu để đánh đổi chất lượng--chi phí~\citep{ong2024routellm}.
\textbf{RAGRouter} nghiên cứu routing trong các hệ retrieval-augmented generation nơi tài liệu truy hồi
ảnh hưởng đến hành vi của mô hình hạ nguồn~\citep{zhang2025ragrouter}. Dự án này thích nghi ý tưởng routing
cho first-stage scientific retrieval: với mỗi query, chọn BM25, Dense, hoặc Hybrid trước khi trả về tài
liệu đã xếp hạng.

\subsection{Fine-tuning hiệu quả, Reranking, QPP và Conformal control}
\textbf{QLoRA} cho phép fine-tune LLM tiết kiệm bộ nhớ bằng 4-bit quantization và LoRA
adapters~\citep{dettmers2023qlora}, khiến Small LLM routing khả thi trên Colab. Router LLM hiện dùng
\textbf{Qwen2.5-0.5B} như một small instruction model~\citep{qwen2024report}. Cho Phase~3, tầng reranking
dùng các cross-encoder nhẹ thuộc họ Sentence-Transformers MS~MARCO~\citep{sbert_msmarco_ce}; thiết kế này
cũng được thúc đẩy bởi BERT-style passage reranking, nơi cross-encoder chấm điểm cặp query--passage một
cách kết hợp và đã chứng minh hiệu quả trên các bài toán ranking kiểu MS~MARCO~\citep{nogueira2019passage}.
\textbf{MiniLM} liên quan vì nó nén transformer thông qua self-attention distillation~\citep{wang2020minilm}.

Hai nền tảng lý thuyết then chốt cho Phase~3 là Query Performance Prediction và Conformal Risk Control.
\textbf{QPP} cho neural retrieval được khảo sát bởi Faggioli et al., những người chỉ ra rằng QPP cho dense
và fused retrieval khó hơn cho lexical retrieval~\citep{faggioli2023qpp}; chúng tôi dùng các predictor cổ
điển như Normalized Query Commitment (NQC)~\citep{shtok2012nqc} và các thống kê score. \textbf{Conformal
Risk Control} (CRC)~\citep{angelopoulos2024crc} và \textbf{Learn-Then-Test}~\citep{angelopoulos2021ltt}
cung cấp các bảo đảm finite-sample về một loss đơn điệu, mà chúng tôi áp dụng để chọn ngưỡng
\emph{khi-nào-rerank} với một bảo đảm về phần nDCG bị mất so với Always-Rerank. Việc áp conformal cho
\emph{quyết định khi nào rerank} (thay vì cho prediction set) là một góc còn ít được khai thác và là nguồn
novelty của báo cáo.
